\hypertarget{appendix-f28eeec97a3d}{}\label{referencestart:appendix-f28eeec97a3d}
\section{Spline bases and curve evaluation}\label{appendix-f28eeec97a3d-spline-bases-and-curve-evaluation}

\textbf{Curriculum depth:} Developed · \textbf{First seen:} University courses · \textbf{Latest development:} University courses

\begin{allmightroledefinition}

\subsection{Definition and mechanism}\label{appendix-f28eeec97a3d-definition-and-mechanism}

A spline is a piecewise polynomial curve whose neighboring pieces join under explicit continuity constraints. A cubic Bézier segment uses four control points and Bernstein weights; the weights sum to one, so the curve lies in their convex hull. A B-spline uses local basis functions so moving one control point changes only nearby spans rather than the whole curve. Knots mark where polynomial pieces meet; repeated knots reduce continuity. Catmull--Rom curves interpolate their control points and can be convenient for motion paths, while Bézier curves expose endpoint tangents through their handles. To evaluate a curve, compute basis weights at parameter \(t\) and take their weighted sum of control points. Position, first derivative, and parameterization speed matter separately: even a smooth path may speed up or slow down unless its parameter is adjusted.

\end{allmightroledefinition}

\subsection{Prerequisites}\label{appendix-f28eeec97a3d-prerequisites}

\begin{itemize}
\tightlist
\item
  \href{../07-Vision-XR-and-Playable-Worlds/main.pdf\#nameddest=appendix-ce0341f7ea0d}{Basis coordinates and dimension (Vision, XR, and Playable Worlds, Reference notes: Basis coordinates and dimension)}
\end{itemize}

\subsection{Used by}\label{appendix-f28eeec97a3d-used-by}

\begin{itemize}
\item
  {Worlds, Images and Inference}
\item
  \href{../07-Vision-XR-and-Playable-Worlds/main.pdf\#nameddest=appendix-ddbec701898c}{Bézier curves and De Casteljau evaluation (Vision, XR, and Playable Worlds, Reference notes: Bézier curves and De Casteljau evaluation)}
\item
  \href{../07-Vision-XR-and-Playable-Worlds/main.pdf\#nameddest=appendix-db3aef8f2fd6}{Cubic Hermite interpolation (Vision, XR, and Playable Worlds, Reference notes: Cubic Hermite interpolation)} \#\# Full teaching explanation
\end{itemize}

\hyperref[03-visual-computing-i-ch001-curves-as-controllable-functions]{Curves as controllable functions} · \href{../05-Spatial-Worlds-Engine-State-and-Simulation/main.pdf\#nameddest=04-visual-computing-ii-ch029-animation-curves-and-motion-paths}{Animation curves and motion paths (Spatial Worlds, Engine State, and Simulation, Ch. 10)}
